% =====================================================================
\section{Data Feasibility Study and Evaluation Plan}\label{sec:exp}
% =====================================================================

\subsection{Data Collection Setup}
We polled the Taichung vehicle API~\cite{dataset83558} every 10 minutes
from 2026-10-05 21:01 to 2026-11-03 19:42 (about 29 days), stored each
distinct response as a snapshot, and deduplicated GPS points by vehicle
and timestamp. We downloaded the collection-point
dataset~\cite{dataset84004} once and joined the two datasets by license
plate. The API returns only the latest position of each vehicle, so
history exists only if it is collected continuously.
Table~\ref{tab:data} summarizes the results.

\begin{table}[htbp]
\caption{Open Data Feasibility Measurements}
\label{tab:data}
\centering
\footnotesize
\begin{tabular}{@{}p{0.56\columnwidth}r@{}}
\toprule
\textbf{Item} & \textbf{Value} \\
\midrule
\multicolumn{2}{@{}l}{\emph{Vehicle feed (29 days)}} \\
Snapshots captured / scheduled          & 4{,}100 / 4{,}169 (98.3\%) \\
Vehicles observed                       & 668 \\
Deduplicated GPS points                 & 980{,}775 \\
Inferred snapshot interval              & $\approx$5 min \\
Vehicles silent $>$30 min (one snapshot) & 118 / 663 (18\%) \\
GPS points with zero speed              & 81\% \\
\midrule
\multicolumn{2}{@{}l}{\emph{Collection points}} \\
Collection points                       & 20{,}090 \\
Curbside / fixed-point                  & 95.1\% / 4.8\% \\
Plates found in vehicle feed            & 358 / 374 (95.7\%) \\
Points covered by matched plates        & 19{,}307 (96.1\%) \\
Points geocoded by local matching       & 19{,}216 (95.6\%) \\
\bottomrule
\end{tabular}
\end{table}

\subsection{Findings}
\emph{Update frequency.} The platform declares a 10-minute
update~\cite{dataset83558}. In 1{,}518 sampled snapshots, however, the
newest GPS fix was consistently about two minutes older than our fetch
time (median 123~s), and a manual check confirmed a snapshot at
five-minute boundaries. We therefore infer that the API publishes a
snapshot roughly every five minutes, far coarser than the 20-second
reports used internally by the city's system~\cite{tcepb2022report}. This
is an inference from limited observation, not an official guarantee.
With stops that last only 5 to 30 minutes~\cite{tcepb2026spring}, a delay
of a few minutes can change whether leaving now is feasible, which is why
data age $\Delta t(s)$ enters the buffer in~(\ref{eq:buffer}).

\emph{Staleness.} In the snapshot of 2026-11-03 19:42, 118 of 663
vehicles had not reported for more than 30 minutes and 8 for more than a
day. Some of these trucks had finished their shifts, so staleness must be
judged against the schedule before being treated as an error.

\emph{Entity matching.} The \texttt{location} field is free text, not a
point identifier. In a random sample of 2{,}000 points, 18.4\% referred to
treatment plants or incinerators, which suggests mid-route unloading that
may delay later stops. License plates, however, use the same format in
both datasets: 95.7\% of the plates in the point dataset appear in the
vehicle feed, covering 96.1\% of the points. Plate matching therefore
links each point to its truck, although not stop by stop.

\emph{Geocoding.} Collection points have addresses but no coordinates.
Matching normalized addresses against the city's 1{,}316{,}674 official
address points located 95.6\% of the points (88.9\% exact matches). The
remaining 4.4\% require an external geocoder or manual review, and the
coordinates have not yet been checked point by point.

\emph{Schedule baseline.} The point dataset lists start and end times for
each weekday, so a schedule-based ETA is available for almost every point
from the start. Only 50 points have no general-garbage slot.

\subsection{Evaluation Plan}
Continuous collection had reached 1{,}710{,}839 deduplicated observations
from 669 vehicles by 2026-11-19 and is ongoing. We will split data by
time into training, validation, and test periods to avoid leakage. The
planned comparisons are:
\begin{itemize}
  \item \emph{ETA:} schedule, distance/speed, and historical-average
        baselines~\cite{rawat2026bus}, compared by MAE, RMSE, and error
        distribution against arrival times inferred from GPS traces.
  \item \emph{Recommendation:} the nearest-point strategy versus the
        slack-based strategy, compared by catch rate, misjudgment rate
        (recommended as feasible but missed), and the success rate of
        alternatives.
  \item \emph{Robustness:} stale data, API failure, and empty candidate
        sets, checking that the agent falls back correctly and that every
        number it reports is traceable to a tool result.
\end{itemize}
All evaluation results are pending. A prerequisite is a reliable way to
infer actual arrival times from five-minute snapshots, which we are
validating first.
